%% ─── Encoding & fonts ────────────────────────────────────────────────────────
\usepackage[T1]{fontenc}
\usepackage[utf8]{inputenc}
\usepackage{lmodern}          % scalable CM fonts (no pixelation in PDF)
\usepackage{microtype}        % optical margin alignment & character protrusion

%% ─── Page layout ─────────────────────────────────────────────────────────────
\usepackage{geometry}
\geometry{left=18mm, right=18mm, top=18mm, bottom=16mm}
\linespread{1.3}

%% ─── Core mathematics ────────────────────────────────────────────────────────
\usepackage{amsmath, amssymb, amsthm, textcomp}
\usepackage{mathtools}        % extends amsmath: \coloneqq, paired delimiters …
\usepackage{bbm}              % \mathbbm{1} for indicator functions
\usepackage{cancel}           % \cancel, \bcancel, \xcancel
\usepackage{esint}            % extended integral signs (\oiint etc.)
\usepackage{stmaryrd}
\allowdisplaybreaks

%% ─── Theorem environments (amsthm counters) ──────────────────────────────────
\newtheorem{theorem}{Theorem}[section]
\renewcommand{\thetheorem}{\arabic{section}.\arabic{theorem}}
\newtheorem{corollary}{Corollary}[theorem]
\newtheorem{lemma}[theorem]{Lemma}
\newtheorem{proposition}[theorem]{Proposition}
\newtheorem{algorithm}{Algorithm}[section]
\newtheorem{method}[theorem]{Method}

\theoremstyle{definition}
\newtheorem{definition}{Definition}[section]
\renewcommand{\thedefinition}{\arabic{section}.\arabic{definition}}
\newtheorem*{example}{Example}
\newtheorem{assumption}{Assumption}[section]
\newtheorem{exercise}{Exercise}[chapter]

\theoremstyle{remark}
\newtheorem*{remark}{Remark}
\newtheorem{review}{Review}

%% ─── Graphics & colour ───────────────────────────────────────────────────────
\usepackage{graphicx}
\usepackage{float}
\usepackage{xcolor}
\usepackage{tikz}

%% ─── Miscellaneous utilities ─────────────────────────────────────────────────
\usepackage{multicol}
\usepackage[normalem]{ulem}
\usepackage{marvosym}
\usepackage{algpseudocode}
\usepackage{url}
\usepackage{csquotes}

%% ─── Code listings ───────────────────────────────────────────────────────────
\usepackage{listings}
\lstset{
  language=Python,
  basicstyle=\ttfamily\small,
  aboveskip={1.0\baselineskip},
  belowskip={1.0\baselineskip},
  columns=fixed,
  extendedchars=true,
  breaklines=true,
  tabsize=4,
  prebreak=\raisebox{0ex}[0ex][0ex]{\ensuremath{\hookleftarrow}},
  frame=lines,
  showtabs=false, showspaces=false, showstringspaces=false,
  keywordstyle=\color[rgb]{0.627,0.126,0.941},
  commentstyle=\color[rgb]{0.133,0.545,0.133},
  stringstyle=\color[rgb]{1,0,0},
  numbers=left, numberstyle=\small, stepnumber=1, numbersep=10pt,
  captionpos=t,
  escapeinside={\%*}{*)}
}

%% ─── Hyperlinks (load last among packages) ───────────────────────────────────
\usepackage[colorlinks=true,
            linkcolor=blue!70!black,
            citecolor=green!55!black,
            urlcolor=blue!80!black]{hyperref}
\usepackage[open,openlevel=1]{bookmark}

%% ─── Section / chapter numbering ─────────────────────────────────────────────
\renewcommand{\thechapter}{\Roman{chapter}}
\renewcommand*\thesection{\arabic{section}}
\providecommand*\theHsection{}
\providecommand*\theHtheorem{}
\providecommand*\theHcorollary{}
\providecommand*\theHalgorithm{}
\providecommand*\theHdefinition{}
\providecommand*\theHassumption{}
\providecommand*\theHexercise{}
\providecommand*\theHreview{}
\renewcommand*\theHsection{\thechapter.\arabic{section}}
\renewcommand*\theHtheorem{\theHsection.\arabic{theorem}}
\renewcommand*\theHcorollary{\theHtheorem.\arabic{corollary}}
\renewcommand*\theHalgorithm{\theHsection.\arabic{algorithm}}
\renewcommand*\theHdefinition{\theHsection.\arabic{definition}}
\renewcommand*\theHassumption{\theHsection.\arabic{assumption}}
\renewcommand*\theHexercise{\thechapter.\arabic{exercise}}
\renewcommand*\theHreview{\arabic{review}}
\DeclareMathAlphabet{\mathcal}{OMS}{cmsy}{m}{n}
\makeatletter
\patchcmd{\l@chapter}{1.5em}{2.4em}{}{}
\makeatother

%% ─── Math operators ──────────────────────────────────────────────────────────
\AtBeginDocument{\let\div\relax  \DeclareMathOperator{\div}{div}}
\AtBeginDocument{\let\ran\relax  \DeclareMathOperator{\ran}{ran}}
\DeclareMathOperator{\dom}{dom}
\DeclareMathOperator{\var}{Var}
\DeclareMathOperator{\cov}{Cov}
\DeclareMathOperator*{\sign}{sign}
\DeclareMathOperator{\tr}{tr}
\DeclareMathOperator{\rank}{rank}
\DeclareMathOperator{\sgn}{sgn}
\DeclareMathOperator{\diam}{diam}
\DeclareMathOperator{\Int}{int}        % interior of a set
\DeclareMathOperator{\co}{co}          % convex hull
\DeclareMathOperator{\grad}{grad}
\DeclareMathOperator{\curl}{curl}
\DeclareMathOperator*{\esssup}{ess\,sup}
\DeclareMathOperator*{\essinf}{ess\,inf}
\DeclareMathOperator*{\spa}{span}
\DeclareMathOperator*{\dist}{dist}
\DeclareMathOperator*{\support}{supp}

%% ─── Number fields ───────────────────────────────────────────────────────────
\newcommand{\R}{\mathbb{R}}
\newcommand{\N}{\mathbb{N}}
\newcommand{\Z}{\mathbb{Z}}
\newcommand{\C}{\mathbb{C}}
\newcommand{\Q}{\mathbb{Q}}
\newcommand{\K}{\mathbb{K}}            % generic field (R or C)
\newcommand{\T}{\mathbb{T}}            % torus
\newcommand{\F}{\mathbb{F}}
\newcommand{\ind}{\mathbbm{1}}         % indicator function \ind_A

%% ─── Paired delimiters (use starred form for auto-size, e.g. \norm*{f}) ──────
\DeclarePairedDelimiter{\abs}{\lvert}{\rvert}
\DeclarePairedDelimiter{\norm}{\lVert}{\rVert}
\DeclarePairedDelimiter{\set}{\{}{\}}
\DeclarePairedDelimiter{\ceil}{\lceil}{\rceil}
\DeclarePairedDelimiter{\floor}{\lfloor}{\rfloor}
\newcommand{\ip}[2]{\langle #1,\, #2 \rangle}       % inner product (fixed size)
\newcommand{\IP}[2]{\left\langle #1,\, #2 \right\rangle} % inner product (auto-size)

%% ─── Greek letter shortcuts ──────────────────────────────────────────────────
\newcommand{\eps}{\varepsilon}
\newcommand{\vp}{\varphi}
\newcommand{\lam}{\lambda}
\newcommand{\Lam}{\Lambda}
\newcommand{\sig}{\sigma}
\newcommand{\om}{\omega}
\newcommand{\Om}{\Omega}
\newcommand{\al}{\alpha}
\newcommand{\be}{\beta}
\newcommand{\ga}{\gamma}
\newcommand{\Ga}{\Gamma}
\newcommand{\de}{\delta}
\newcommand{\De}{\Delta}
\newcommand{\ka}{\kappa}
\newcommand{\tht}{\theta}
\newcommand{\Tht}{\Theta}

%% ─── Calculus & analysis ─────────────────────────────────────────────────────
\newcommand{\del}{\partial}                          % partial derivative / boundary
\newcommand{\nab}{\nabla}                            % nabla
\newcommand{\Lap}{\Delta}                            % Laplacian
\newcommand{\dd}{\mathop{}\!\mathrm{d}}              % differential: \int f \dd x
\newcommand{\ddt}{\tfrac{\mathrm{d}}{\mathrm{d}t}}  % d/dt

%% ─── Topology & convergence ──────────────────────────────────────────────────
\newcommand{\cl}[1]{\overline{#1}}                   % closure: \cl{A}
\newcommand{\wto}{\rightharpoonup}                   % weak convergence
\newcommand{\wsto}{\overset{*}{\rightharpoonup}}     % weak-* convergence
\newcommand{\into}{\hookrightarrow}                  % injection
\newcommand{\onto}{\twoheadrightarrow}               % surjection
\newcommand{\xto}[1]{\xrightarrow{#1}}               % labelled arrow: \xto{L^p}

%% ─── Miscellaneous math ──────────────────────────────────────────────────────
\newcommand{\id}{\mathrm{id}}                        % identity operator
\newcommand{\st}{\;\text{s.t.}\;}                    % such that
\newcommand{\cco}{\overline{\co}}                    % closed convex hull

%% ─── Page style ──────────────────────────────────────────────────────────────
\usepackage{fancyhdr}
\setlength{\headheight}{13.59999pt}
\addtolength{\topmargin}{-1.59999pt}
\pagestyle{fancy}
\lhead{}\chead{}\rhead{}
\lfoot{}\cfoot{}\rfoot{\thepage}
\renewcommand{\headrulewidth}{0pt}
\renewcommand{\footrulewidth}{0pt}

%% ─── Title style ─────────────────────────────────────────────────────────────
\newcommand{\linia}{\rule{\linewidth}{0.5pt}}
\makeatletter
\renewcommand{\maketitle}{%
  \begin{center}
  \vspace{2ex}
  {\huge \textsc{\@title}}
  \vspace{1ex}\\
  \linia\\
  \@author \hfill \@date
  \vspace{4ex}
  \end{center}}
\makeatother
